\textbf{Phần 1. Trắc nghiệm lựa chọn}
\setcounterpageref{ex}{0}

\begin{ex} %[4D8N38]
Kết quả của phép nhân $23\,142 \times 3$ là
\choice
{$69\,326$}
{$69\,416$}
{$68\,426$}
{\True $69\,426$}
\loigiai{
$23\,142 \times 3 = 69\,426$.\\
\textbf{Đáp án của câu hỏi là D}
}
\end{ex}

\begin{ex} %[4D8N38]
Kết quả của phép tính $120\,000 \times 4$ là
\choice
{\True $480\,000$}
{$460\,000$}
{$400\,000$}
{$500\,000$}
\loigiai{
$120\,000 \times 4 = 480\,000$.\\
\textbf{Đáp án của câu hỏi là A}
}
\end{ex}

\begin{ex} %[4D8H38]
Tìm $x$ biết: $x : 6 = 14\,250$.
\choice
{$84\,500$}
{$85\,200$}
{$86\,500$}
{\True $85\,500$}
\loigiai{
$x = 14\,250 \times 6 = 85\,500$.\\
\textbf{Đáp án của câu hỏi là D}
}
\end{ex}

\begin{ex} %[4D8H38]
Kết quả của phép tính $45\,678 \times 5$ là
\choice
{$228\,380$}
{\True $228\,390$}
{$225\,390$}
{$228\,490$}
\loigiai{
$45\,678 \times 5 = 228\,390$.\\
\textbf{Đáp án của câu hỏi là B}
}
\end{ex}

\begin{ex} %[4D8H38]
Một hộp có $125$ chiếc bút. Hỏi $8$ hộp như thế có tất cả bao nhiêu chiếc bút?
\choice
{$900$ chiếc}
{$1100$ chiếc}
{\True $1000$ chiếc}
{$800$ chiếc}
\loigiai{
$125 \times 8 = 1000$ chiếc.\\
\textbf{Đáp án của câu hỏi là C}
}
\end{ex}

\begin{ex} %[4D8H38]
Giá trị của biểu thức $102\,300 \times 3 + 45\,000$ là
\choice
{\True $351\,900$}
{$341\,900$}
{$350\,900$}
{$361\,900$}
\loigiai{
$306\,900 + 45\,000 = 351\,900$.\\
\textbf{Đáp án của câu hỏi là A}
}
\end{ex}

\begin{ex} %[4D8V38]
Một phân xưởng may mỗi ngày may được $1250$ bộ quần áo. Hỏi trong một tuần lễ làm việc 6 ngày, phân xưởng đó may được bao nhiêu bộ quần áo?
\choice
{$7200$ bộ}
{\True $7500$ bộ}
{$7800$ bộ}
{$8000$ bộ}
\loigiai{
$1250 \times 6 = 7500$ bộ.\\
\textbf{Đáp án của câu hỏi là B}
}
\end{ex}

\begin{ex} %[4D8V38]
Một ô tô chở hàng nặng $3$ tấn $500\text{ kg}$. Hỏi 4 ô tô chở như thế chở được tất cả bao nhiêu tấn hàng?
\choice
{$12\text{ tấn}$}
{$15\text{ tấn}$}
{\True $14\text{ tấn}$}
{$16\text{ tấn}$}
\loigiai{
$3500 \times 4 = 14\,000\text{ kg} = 14\text{ tấn}$.\\
\textbf{Đáp án của câu hỏi là C}
}
\end{ex}

\begin{ex} %[4D8V38]
Tính giá trị của biểu thức: $25\,000 \times 4 - 30\,000 : 3$.
\choice
{$80\,000$}
{$100\,000$}
{$70\,000$}
{\True $90\,000$}
\loigiai{
$100\,000 - 10\,000 = 90\,000$.\\
\textbf{Đáp án của câu hỏi là D}
}
\end{ex}

\begin{ex} %[4D8V38]
Một trường tiểu học mua 8 bộ bàn ghế cho phòng họp, mỗi bộ có giá $1\,450\,000$ đồng. Tổng số tiền mua bàn ghế là
\choice
{$11\,400\,000$ đồng}
{\True $11\,600\,000$ đồng}
{$11\,800\,000$ đồng}
{$12\,000\,000$ đồng}
\loigiai{
$1\,450\,000 \times 8 = 11\,600\,000$ đồng.\\
\textbf{Đáp án của câu hỏi là B}
}
\end{ex}

\begin{ex} %[4D8V38]
Tìm chữ số $a$ biết: $12a,400 \times 3 = 373\,200$.
\choice
{\True $a = 4$}
{$a = 3$}
{$a = 5$}
{$a = 2$}
\loigiai{
$373\,200 : 3 = 124\,400 \Rightarrow a = 4$.\\
\textbf{Đáp án của câu hỏi là A}
}
\end{ex}

\begin{ex} %[4D8C38]
Khi nhân một số có 5 chữ số với 4, một học sinh đã viết nhầm chữ số hàng trăm của số đó từ 3 thành 8 nên tích tăng thêm bao nhiêu đơn vị?
\choice
{$500$ đơn vị}
{$1500$ đơn vị}
{\True $2000$ đơn vị}
{$5000$ đơn vị}
\loigiai{
Thừa số thứ nhất tăng thêm: $800 - 300 = 500$.\\
Tích tăng thêm: $500 \times 4 = 2000$ đơn vị.\\
\textbf{Đáp án của câu hỏi là C}
}
\end{ex}

\textbf{Phần 2. Trắc nghiệm đúng sai}
\setcounterpageref{ex}{0}

\begin{ex} %[4D8H38]
Xét tính đúng sai của các phép nhân:
\choiceTF
{\True $12\,345 \times 2 = 24\,690$.}
{$25\,000 \times 4 = 1\,000\,000$.}
{\True $25\,000 \times 4 = 100\,000$.}
{\True $150\,000 \times 6 = 900\,000$.}
\loigiai{
\begin{itemize}
    \item $12\,345 \times 2 = 24\,690$. Đúng.
    \item $25\,000 \times 4 = 100\,000 \neq 1\,000\,000$. Ý 2 Sai, Ý 3 Đúng.
    \item $150\,000 \times 6 = 900\,000$. Đúng.
\end{itemize}
\begin{itemize}
    \item a) \textbf{Đúng}.
    \item b) \textbf{Sai}.
    \item c) \textbf{Đúng}.
    \item d) \textbf{Đúng}.
\end{itemize}
}
\end{ex}

\begin{ex} %[4D8V38]
Một cửa hàng nhập về 5 bao đường, mỗi bao nặng $50\text{ kg}$. Giá mỗi ki-lô-gam đường là $22\,000$ đồng. Xét tính đúng sai:
\choiceTF
{\True Cửa hàng đã nhập về tất cả $250\text{ kg}$ đường.}
{\True Số đường đó tương đương $2{,}5$ tạ đường.}
{Tổng số tiền mua đường là $5\,000\,000$ đồng.}
{\True Tổng số tiền mua đường là $5\,500\,000$ đồng.}
\loigiai{
\begin{itemize}
    \item $5 \times 50 = 250\text{ kg}$. Đúng.
    \item $250\text{ kg} = 2{,}5$ tạ. Đúng.
    \item Tiền: $250 \times 22\,000 = 5\,500\,000$ đồng $\neq 5\,000\,000$ đồng. Ý 3 Sai, Ý 4 Đúng.
\end{itemize}
\begin{itemize}
    \item a) \textbf{Đúng}.
    \item b) \textbf{Đúng}.
    \item c) \textbf{Sai}.
    \item d) \textbf{Đúng}.
\end{itemize}
}
\end{ex}

\textbf{Phần 3. Trả lời ngắn}
\setcounterpageref{ex}{0}

\begin{ex} %[4D8H38]
Tính tích: $15\,000 \times 6$.
\shortans{90000}
\loigiai{
\textbf{Đáp số:} $90000$\\
$90\,000$.
}
\end{ex}

\begin{ex} %[4D8H38]
Tính tích: $125\,000 \times 8$.
\shortans{1000000}
\loigiai{
\textbf{Đáp số:} $1000000$\\
$1\,000\,000$.
}
\end{ex}

\begin{ex} %[4D8V38]
Tìm $x$ biết: $x : 7 = 35\,000$.
\shortans{245000}
\loigiai{
\textbf{Đáp số:} $245000$\\
$x = 35\,000 \times 7 = 245\,000$.
}
\end{ex}

\begin{ex} %[4D8V38]
Một ngày có $24$ giờ, mỗi giờ có $60$ phút. Tính số phút trong 3 ngày.
\shortans{4320}
\loigiai{
\textbf{Đáp số:} $4320$\\
$24 \times 60 \times 3 = 4320$ phút.
}
\end{ex}

\begin{ex} %[4D8V38]
Một cửa hàng bán được 8 thùng dầu, mỗi thùng chứa $25\text{ l}$ dầu. Giá mỗi lít dầu là $20\,000$ đồng. Tính tổng số tiền thu được theo nghìn đồng (chỉ ghi số nghìn đồng).
\shortans{4000}
\loigiai{
\textbf{Đáp số:} $4000$\\
$8 \times 25 \times 20\,000 = 4\,000\,000$ đồng $= 4000$ nghìn đồng.
}
\end{ex}

\begin{ex} %[4D8C38]
Tích của hai số là $45\,000$. Nếu một thừa số gấp lên 3 lần và thừa số kia giữ nguyên thì tích mới là bao nhiêu?
\shortans{135000}
\loigiai{
\textbf{Đáp số:} $135000$\\
$45\,000 \times 3 = 135\,000$.
}
\end{ex}
